\begin{center}\bfseries 25.\end{center}
\noindent\hspace*{\parindent}Aus etwa derselben Zeit d. h. also aus dem letzten Drittel des Jahres 1808 (s. W. II, S. 398) stammen andere handschriftliche Aufzeichnungen, welche allgemeiner die Theorie der aus einer dritten Einheitswurzel \(\rho\) gebildeten komplexen ganzen Zahlen zum Gegenstande haben (W. II, S. 387—398).

Ist \(a+b\rho+c\rho^2\), wo \(\rho\) eine dritte Einheitswurzel bedeutet, eine komplexe Zahl mit rationalen \(a,b,c\), so lässt sich eine nächstgelegene ganze komplexe
\newpage\SourceLeaf{68}\ArithmeticHead{61}{R}{KUBISCHE UND BIQUADRATISCHE RESTE. KOMPLEXE GANZE ZAHLEN.}
\noindent Zahl so bestimmen, dass der »Determinant« (d. i. die Norm) der Differenz kleiner als 1 wird. Infolge hiervon ist der \textsc{Euklidische} Algorithmus anwendbar und folglich herrscht auch im Gebiete dieser komplexen ganzen Zahlen eindeutige Zerlegbarkeit. Hiermit hatte \textsc{Gauss} die Grundlage gewonnen, um das »letzte \textsc{Fermat}sche Theorem« für den dritten Grad, nämlich die Unlösbarkeit der kubischen Gleichung
\[
x^3+y^3+z^3=0
\]
nicht nur in rationalen, sondern allgemeiner in komplexen ganzen Zahlen der gedachten Art zu beweisen. Dieselben Identitäten benutzend, deren auch \textsc{Legendre} und die späteren Forscher sich bedient haben, zeigt \textsc{Gauss} zunächst, dass, wenn jene kubische Gleichung bestehen soll, eine der drei Zahlen \(x,y,z\) durch \(1-\rho\) teilbar sein muss, und leitet dann (auf zweifachem Wege) aus einer vorausgesetzten Lösung in teilerfremden Zahlen eine zweite (und dritte) von gleicher Beschaffenheit her. Bei der auf dem zweiten Wege erhaltenen neuen Lösung ist die durch \(1-\rho\) teilbare Zahl nicht mehr ebenso oft durch \(1-\rho\) teilbar wie früher, wodurch der Beweis erbracht ist. Hier tritt auch schon der später von \textsc{Kummer}\footnote{\textit{De numeris complexis, qui radicibus unitatis et numeris integris realibus constant}, Vratislaviae typis universitatis, 1844; vergl. \textsc{Crelles} Journal für Mathematik \textbf{35} (1847), S. 362.}\ begründete Satz auf, dass eine komplexe ganze Zahl \(f(\rho)\), deren analytischer Modul gleich 1 ist, \(\pm\rho^n\) sein muss.

Was endlich die Fragmente W. VIII, S. 21 anbelangt, so finden sich die beiden ersten von ihnen in \textsc{Gauss}’ Handbuch 18, Bd vom Oktober 1805, S. 201—204, also nach den schon bei dem Fragmente W. VIII, S. 15 erwähnten Beobachtungen vom August 1808, aber vor diesem Fragmente selbst und sind demnach zwischen denselben Daten verfasst wie dieses. Die im zweiten derselben behandelte kubische Form
\[
x^3+ny^3+n^2z^3-3nxyz
\]
tritt schon auf S. 151 jenes Handbuchs auf d. i. vor dem Beweise für die Irreduzibilität der allgemeinen Kreisteilungsgleichung, den man auf den 12. Juni 1808 datieren darf, und findet sich auch an andern Stellen des Nachlasses, so z. B. auf einer leider undatierten Rechnung für die Sternwarte in Göttingen in Ec 1, woher auch das dritte der Fragmente (W. VIII, S. 23)
\newpage\SourceLeaf{69}\ArithmeticHead{62}{L}{BACHMANN, ÜBER GAUSS’ ZAHLENTHEORETISCHE ARBEITEN.}
\noindent entnommen ist. Alle diese Notizen stammen wohl ziemlich aus der gleichen Zeit. Die gedachten Fragmente lassen nun ersehen, dass \textsc{Gauss} auch die aus einer beliebigen Kubikwurzel \(r=\sqrt[3]{n}\) für ein ganzzahliges \(n\) gebildeten komplexen ganzen Zahlen in den Bereich seiner Untersuchungen gezogen hat. Ist \(a+b\nu+c\nu^2\) eine solche, so ist ihre Norm
\[
N=a^3+nb^3+n^2c^3-3nabc
\]
und \(n\) ist kubischer Rest in Bezug auf jeden Teiler von \(N\). Durch die Substitution
\[
\begin{aligned}
x&=ax'+ncy'+nbz'\\
y&=bx'+ay'+ncz'\\
z&=cx'+by'+az'
\end{aligned}
\]
geht die Form \(x^3+ny^3+n^2z^3-3nxyz\) über in
\[
N\mathbin{.}(x'^3+ny'^3+n^2z'^3-3nx'y'z'),
\]
ein Satz, welcher die Reproduktion dieser Form durch Zusammensetzung mit sich selbst ausspricht und eine Grundlage bildet für die Auflösung der Gleichung
\[
x^3+ny^3+n^2z^3-3nxyz=1,
\]
mit der sich \textsc{Gauss} (Nr. 137 T.) vom 23. Dezember 1808 an ebenfalls beschäftigt hat.
